\documentclass[10pt,a4paper]{amsart}
\usepackage[margin=19mm]{geometry}
\usepackage{amsmath,amssymb,mathtools}
\usepackage[scr=rsfs,cal=euler]{mathalfa}
\usepackage{xfrac}
\usepackage[unicode,hidelinks,bookmarks=false]{hyperref}
\newcommand{\ie}{i.e.\ }
\newcommand{\cf}{cf.\ }
\newcommand{\resp}{respectively\ }
\newcommand{\Subsection}{\subsection}
\providecommand{\nobreakdash}{\nobreak\hbox{-}}
\setlength{\emergencystretch}{2em}
\multlinegap0pt
\usepackage{bm}
\usepackage{bbm}
\usepackage{dsfont}
\usepackage{xspace}
\usepackage{cancel}
\usepackage{mathtools}
\usepackage{amsthm}
\makeatletter
\@ifundefined{theorem}{\newtheorem{theorem}{Theorem}}{}
\@ifundefined{lemma}{\newtheorem{lemma}[theorem]{Lemma}}{}
\makeatother
\title[Abel-Jacobi Map and Symplectic Topology]{Abel-Jacobi Map and Symplectic Topology}
\author{Stephane Tchuiaga}
\date{Source: arXiv:2608.07787v1 (CC BY 4.0)}
\begin{document}
\maketitle

\noindent\emph{Source and licence.} Abel-Jacobi Map and Symplectic Topology. Version arXiv:2608.07787v1 (CC BY 4.0), distributed under the Creative Commons Attribution 4.0 International licence, as recorded in the arXiv licence metadata for this exact version and shown on the abstract page https://arxiv.org/abs/2608.07787v1. Full rights notice: licenses/arxiv-2608.07787-CC-BY-4.0.txt.\par\smallskip

\noindent\textit{Editorial scope.} Counted unit: the Lemma whose source label is \texttt{lem:stability\_lifts} in the article (source0.tex lines 1276--1286, source label \texttt{lem:stability\_lifts}), delivered together with the article's own complete original proof of it (source0.tex lines 1288--1365, one \texttt{proof} environment of 78 lines). No mathematics is written, repaired or re-derived: statement and proof are the article's own text line for line. Required context retained uncounted: none. Every definition, display and label of those lines is retained unchanged. Omitted from this package and not redistributed: 1--1275, 1287--1287, 1366--2139. Nothing inside a retained excerpt is omitted or altered: every retained body is delivered exactly as the article's lines are written, so no omission receipt of any kind accompanies this package. Citations: the retained excerpts carry no citation command at all, so no bibliography entry of the article is transcribed and no reference is redistributed. Reference adaptation: none was needed and none was made; every reference inside the delivered unit resolves inside it, so no unresolved reference or citation is produced. Printed slips kept literal: no printed word, symbol, hypothesis or number of the counted statement or of its proof was altered. Layout and class: the article's own class is \texttt{article} with packages including inputenc, fontenc, amsmath, amssymb, amsthm, hyperref, tikz, amscd, xurl, etoolbox, url; the wrapper is the lane's pinned \texttt{amsart} editorial wrapper at 10pt on A4, which restates the theorem-like environments the retained text needs and adds the article's own macro definitions that the retained text needs (none), with extra wrapper packages (bm, bbm, dsfont, xspace, cancel, mathtools). The delivered numbering is the unit's own. Assets: no retained span contains a figure environment, a graphics-inclusion command, a PDF-page inclusion, a table or a TikZ picture; no image file is copied into this package, the package declares no asset dependency and no image file is redistributed with it. Licence: version arXiv:2608.07787v1 of this article is distributed under the Creative Commons Attribution 4.0 International licence, as recorded in the arXiv licence metadata for this exact version and shown on the abstract page https://arxiv.org/abs/2608.07787v1. A full rights notice is delivered at licenses/arxiv-2608.07787-CC-BY-4.0.txt. Characters: every retained excerpt is pure ASCII, so the delivered engine, XeLaTeX with its default Latin Modern text font, reports no missing character.
	\begin{lemma}[Local Stability of Lifts]
		\label{lem:stability_lifts}
		Let $\pi:\widetilde{M}\to M$ be a regular covering of a compact Riemannian
		manifold $M$ equipped with the pullback metric. Let $K\subset\widetilde M$ be a
		compact connected subset, and let $\tilde{x}_0\in K$ be a fixed basepoint.
		Suppose that $f_n,f:M\to M$ are continuous maps satisfying $f_n\to f$ uniformly.
		Let $\widetilde f_n,\widetilde f:\widetilde M\to\widetilde M$ be lifts of
		$f_n\circ\pi$ and $f\circ\pi$, respectively, satisfying $\widetilde
		f_n(\tilde{x}_0)\longrightarrow \widetilde f(\tilde{x}_0).$ Then $\widetilde
		f_n|_K\longrightarrow \widetilde f|_K$ uniformly on $K$.
	\end{lemma}

	\begin{proof}
		Since $K$ is compact and $\widetilde f$ is continuous, the set $K_1:=\widetilde
		f(K)$ is a compact subset of $\widetilde M$. Because $\pi$ is a covering local
		isometry, every point $\tilde y\in K_1$ admits a radius $r_{\tilde y}>0$ such
		that the ball $B_{\widetilde M}(\tilde y,r_{\tilde y})$ is contained in an
		evenly covered neighborhood $\widetilde U_{\tilde y}\subset\widetilde M$ and is
		geodesically convex. The family $\left\{ B_{\widetilde M}(\tilde y,r_{\tilde
			y}/2) \right\}_{\tilde y\in K_1}$ is an open cover of the compact set $K_1$.
		Hence there exist finitely many points $\tilde y_1,\ldots,\tilde y_m\in K_1$
		such that
		\[
		K_1\subset \bigcup_{j=1}^{m} B_{\widetilde M} (\tilde y_j,r_j/2),
		\]
		where $r_j:=r_{\tilde y_j}.$ Define
		\[
		\delta_K = \frac14\min_{1\le j\le m}r_j.
		\]
		Then for every $\tilde y\in K_1$, there exists an index $j$ such that $\tilde
		y\in B_{\widetilde M}(\tilde y_j,r_j/2).$ Therefore, $B_{\widetilde M}(\tilde
		y,\delta_K) \subset B_{\widetilde M}(\tilde y_j,r_j) \subset \widetilde
		U_{\tilde y_j}.$ Hence every $\delta_K$-ball centered at a point of $K_1$ is
		contained in a geodesically convex evenly covered neighborhood. Since $f_n\longrightarrow
		f$ uniformly on $M$, there exists $N$ such that for every $n\ge N$,
		\[
		\sup_{x\in M} d_M(f_n(x),f(x)) < \delta_K.
		\]
		Fix $n\ge N$ and $\tilde x\in K$. Put $\tilde y=\widetilde f(\tilde x)\in K_1.$
		The points $f(\pi(\tilde x))$ and $f_n(\pi(\tilde x))$ are at distance less than
		$\delta_K$. Hence there exists a unique minimizing geodesic $\gamma_{n,\tilde
			x}:[0,1]\to M$ joining them. Lift this geodesic to a path $\widetilde\gamma_{n,\tilde
			x}$ starting at $\widetilde\gamma_{n,\tilde x}(0) = \widetilde f(\tilde x).$
		Let $\widetilde g_n(\tilde x) = \widetilde\gamma_{n,\tilde x}(1).$ By
		construction, $\pi(\widetilde g_n(\tilde x)) = f_n(\pi(\tilde x)),$ so
		$\widetilde g_n$ is a lift of $f_n\circ\pi|_K.$ We now estimate the distance
		between the endpoints. Since the length of the lifted geodesic equals the length
		of the original geodesic, we have
		\[
		\operatorname{Length} (\widetilde\gamma_{n,\tilde x}) =
		d_M(f_n(\pi(\tilde x)),f(\pi(\tilde x))) < \delta_K.
		\]
		Because $B_{\widetilde M}(\widetilde f(\tilde x),\delta_K)$ is contained in a
		geodesically convex evenly covered neighborhood, the lifted geodesic remains
		inside this neighborhood. On this neighborhood, $\pi$ is an isometry. Therefore,
		\[
		d_{\widetilde M} (\widetilde g_n(\tilde x),\widetilde f(\tilde x)) =
		d_M(f_n(\pi(\tilde x)),f(\pi(\tilde x))) \le \|f_n-f\|_{C^0}.
		\]
		Consequently,
		\[
		\sup_{\tilde x\in K} d_{\widetilde M} (\widetilde g_n(\tilde x),\widetilde
		f(\tilde x)) \le \|f_n-f\|_{C^0}.
		\tag{a}
		\]
		
		It remains to identify $\widetilde g_n$ with the prescribed lift $\widetilde
		f_n$. By assumption, $\widetilde f_n(\tilde x_0) \longrightarrow \widetilde
		f(\tilde x_0).$ Moreover, applying (a) at $\tilde x_0$ gives $\widetilde
		g_n(\tilde x_0) \longrightarrow \widetilde f(\tilde x_0).$ Hence for all
		sufficiently large $n$, both points $\widetilde f_n(\tilde x_0)$ and
		$\widetilde g_n(\tilde x_0)$ belong to the same evenly covered neighborhood of
		$\widetilde f(\tilde x_0)$. They are lifts of the same point:
		\[
		\pi(\widetilde f_n(\tilde x_0)) = f_n(\pi(\tilde x_0)) =
		\pi(\widetilde g_n(\tilde x_0)).
		\]
		Since $\pi$ is injective on this evenly covered neighborhood, we conclude that
		$\widetilde f_n(\tilde x_0) = \widetilde g_n(\tilde x_0).$ Therefore,
		$\widetilde f_n|_K$ and $\widetilde g_n$ are two lifts of the same map
		$f_n\circ\pi|_K$ which agree at the point $\tilde x_0$. Since $K$ is connected,
		the uniqueness of lifts implies $\widetilde f_n|_K = \widetilde g_n.$ Combining
		this equality with (a), we obtain
		\[
		\sup_{\tilde x\in K} d_{\widetilde M} (\widetilde f_n(\tilde x),\widetilde
		f(\tilde x)) \le \|f_n-f\|_{C^0}.
		\]
		Because $\|f_n-f\|_{C^0}\longrightarrow0,$ we conclude that $\widetilde f_n|_K
		\longrightarrow \widetilde f|_K$ uniformly on $K$.
	\end{proof}

\end{document}
